\documentclass[11pt]{article}
\usepackage[margin=1in]{geometry}
\usepackage{amsmath,amssymb,mathtools}
\usepackage{fontspec,unicode-math}
\setmainfont{Latin Modern Roman}
\setsansfont{Latin Modern Sans}
\setmathfont{Latin Modern Math}
\usepackage{xurl,hyperref}
\hypersetup{colorlinks=true,urlcolor=blue,linkcolor=blue}
\newcommand{\sref}[2]{\hyperlink{source:#1:#2}{[S#2]}}
\newcommand{\cataloguescope}[1]{\paragraph{Recorded mathematical question.}#1}
\begin{document}
\section{Ends of infinite clusters in the nonuniqueness phase}\label{statement}

\subsection{Definitions and mathematical statement}

\paragraph{Shared notation.}$\mathbb N=\{1,2,\ldots\}$, and $\mathbb Z,\mathbb Q,\mathbb R,\mathbb C$ denote the integers, rationals, reals and complex numbers. Any other conventions are specified locally.


Let $G$ be an infinite connected locally finite quasi-transitive graph: its automorphism group has finitely many vertex orbits. In Bernoulli site percolation, vertices are independently open with probability $p$ and clusters are components induced by open vertices. The bond variant instead independently opens edges. An end of an infinite connected locally finite graph is an equivalence class of infinite simple rays, two rays being equivalent when for every finite set of vertices their tails lie in the same component after deleting that set. The cardinal $2^{\aleph_0}$ is the cardinality of the continuum.
\paragraph{Statement recorded in the supplied source.}
For $p\in(0,1)$, if there is more than one infinite cluster almost surely, prove that every infinite cluster has exactly $2^{\aleph_0}$ ends almost surely. State the same conclusion for Bernoulli bond percolation. \sref{statement}{2}
\subsection{Short English statement}
In the nonunique phase on a quasi-transitive graph, must every infinite cluster have continuum many ends?
\subsection{Sources}
\begin{description}
\item[\hypertarget{source:statement:1}{S1}] ULAM.AI and the UnsolvedMath Contributors, \emph{UnsolvedMath: A Curated Collection of Open Mathematics Problems}, catalogue entry. CC BY 4.0. \url{https://huggingface.co/datasets/ulamai/UnsolvedMath}.
\item[\hypertarget{source:statement:2}{S2}] Itai Benjamini and Oded Schramm, \emph{Percolation beyond $\mathbb Z^d$, many questions and a few answers}, Electronic Communications in Probability 1 (1996), 71--82, Section 2, conventions, p. 72; Section 4, Conjecture 5, p. 76. \url{https://emis.de/ft/41630}.
\end{description}
\label{end:statement}
\end{document}
